\chapter{Egzotyczna czterorozmaitość Parka: konstrukcja i topologia}
\label{ch:tg-park}

Dotychczasowy aparat pozwala połączyć trzy pozornie odległe
konstrukcje: grafy plumbingowe, wymierną kontrakcję i teorię
czterorozmaitości symplektycznych. Celem rozdziału jest
\emph{skonstruowanie i zbadanie topologicznie} rozmaitości
$X_7$, której topologicznym modelem jest
$\mathbb{CP}^{2}\#7\overline{\mathbb{CP}}^{\,2}$.
Rozróżnienie jej struktury gładkiej od standardowej wymaga
niezmienników Seiberga--Wittena oraz wyboru komory; ten
ostatni argument przedstawimy w następnym rozdziale.

Wszystkie rozmaitości w tym rozdziale są zwarte i gładkie,
chyba że wyraźnie zaznaczono kategorię topologiczną.
Zakładamy orientacje zgodne z zespoloną orientacją
$\mathbb{CP}^{2}$ i standardową orientacją plumbingu.
Piszemy $\overline{\mathbb{CP}}^{\,2}$ dla przestrzeni
rzutowej z orientacją przeciwną.

\section{Racjonalna powierzchnia eliptyczna jako punkt wyjścia}
\label{sec:tg-park-e1}

\begin{definicja}[Powierzchnia eliptyczna $E(1)$]
Przez $E(1)$ oznaczamy rozdmuchanie dziewięciu punktów
bazowych odpowiednio wybranego pęku krzywych sześciennych
na $\mathbb{CP}^{2}$. W rezultacie dostajemy powierzchnię
zespoloną z odwzorowaniem na $\mathbb{CP}^{1}$, którego
włókno ogólne jest gładką krzywą eliptyczną. Topologicznie
$E(1)\cong\mathbb{CP}^{2}\#9\overline{\mathbb{CP}}^{\,2}$.
\end{definicja}

Wybieramy bazę $h,e_1,\ldots,e_9$ w drugiej homologii,
przy czym $h$ jest klasą prostej w $\mathbb{CP}^{2}$,
a $e_j$ klasą sfery wyjątkowej. Forma przecięcia jest
diagonalna: $h^2=1$, $e_j^2=-1$, $h\cdot e_j=0$
i $e_i\cdot e_j=0$ dla $i\ne j$.
Klasa ogólnego włókna eliptycznego wynosi
$f=3h-e_1-\cdots-e_9$.

\begin{lemat}
\label{lem:tg-park-fiber}
W opisanej bazie zachodzą $f^2=0$ oraz $f\cdot e_j=1$
dla $1\le j\le9$. Po czterech dalszych rozdmuchaniach
otrzymujemy $R:=E(1)\#4\overline{\mathbb{CP}}^{\,2}$,
o niezmiennikach
\[
e(R)=16,\qquad \sigma(R)=-12,\qquad
b_2^+(R)=1,\quad b_2^-(R)=13.
\]
\end{lemat}

\begin{proof}
Z postaci diagonalnej otrzymujemy
$f^2=9-\sum_{j=1}^{9}1=0$,
a $f\cdot e_j=-e_j\cdot e_j=1$.
Każde rozdmuchanie zwiększa charakterystykę Eulera
o jeden, zmniejsza sygnaturę o jeden i dodaje jedną
ujemną składową do formy przecięcia. Ponieważ
$e(\mathbb{CP}^{2})=3$ i
$\sigma(\mathbb{CP}^{2})=1$, dla $R$ mamy
$e(R)=3+13=16$ oraz $\sigma(R)=1-13=-12$.
Równości dla $b_2^\pm$ wynikają z diagonalizacji formy.
\end{proof}

\begin{przyklad}[Klasa sfery po czterech rozdmuchaniach]
\label{ex:tg-park-S}
Niech $e_{10},\ldots,e_{13}$ oznaczają klasy nowych
sfer wyjątkowych. Rozważmy klasę
\[
s:=4f+e_9-2(e_{10}+e_{11}+e_{12}+e_{13})\in H_2(R;\mathbb Z).
\]
Korzystając z $f^2=0$, $f\cdot e_9=1$,
$e_9^2=-1$ oraz z ortogonalności $f$ i $e_9$ do nowych klas
wyjątkowych, obliczamy
\[
s^2=8(f\cdot e_9)+e_9^2
   +4\sum_{j=10}^{13}e_j^2=8-1-16=-9.
\]
Jest to rachunek homologiczny, nie dowód istnienia
gładko zanurzonej sfery reprezentującej $s$.
Konstrukcja geometryczna takiej sfery korzysta z włókien
osobliwych pęku eliptycznego i rozwiązywania dodatnich
przecięć, jak wyjaśniamy niżej.
\end{przyklad}

\begin{center}
\begin{tikzpicture}[>=Stealth,line join=round,line cap=round]
  \node[draw=tgBlue,rounded corners=5pt,fill=tgBlue!9,
        minimum width=2.65cm,minimum height=1cm,align=center,text width=2.3cm,font=\small] (a)
        at (-4.6,0) {$\mathbb{CP}^{2}$\\pęk sześciennych};
  \node[draw=tgTeal,rounded corners=5pt,fill=tgTeal!10,
        minimum width=2.65cm,minimum height=1cm,align=center,text width=2.3cm,font=\small] (b)
        at (0,0) {$E(1)$\\włókna eliptyczne};
  \node[draw=tgPurple,rounded corners=5pt,fill=tgPurple!11,
        minimum width=2.65cm,minimum height=1cm,align=center,text width=2.3cm,font=\small] (c)
        at (4.6,0) {$R$\\$4$ nowe rozdmuchania};
  \draw[-{Stealth[length=2.8mm]},line width=1.65pt,tgTeal]
        (a)--(b);
  \draw[-{Stealth[length=2.8mm]},line width=1.65pt,tgPurple]
        (b)--(c);
\end{tikzpicture}

{\footnotesize Pierwsza strzałka: $9$ rozdmuchań punktów bazowych;
druga: $4$ rozdmuchania punktów podwójnych włókien osobliwych.
$E(1)=\mathbb{CP}^{2}\#9\overline{\mathbb{CP}}^{\,2}$ oraz
$R=E(1)\#4\overline{\mathbb{CP}}^{\,2}$.}
\end{center}

\section{Sześć sfer i graf plumbingowy \texorpdfstring{$C_7$}{C7}}
\label{sec:tg-park-c7}

\begin{definicja}[Konfiguracja $C_7$]
Niech $\Gamma_7$ będzie grafem łańcuchowym o sześciu
wierzchołkach, z wagami $(-9,-2,-2,-2,-2,-2)$.
Konfiguracją $C_7$ nazywamy zwartą czterorozmaitość
uzyskaną przez plumbing wiązek dyskowych zgodny z tym grafem.
Przez $S_0,\ldots,S_5$ oznaczamy sfery zerowe wiązek
dyskowych. Ich klasy tworzą bazę $H_2(C_7;\mathbb Z)$
i spełniają
\[
S_0^2=-9,\quad S_j^2=-2\ (1\le j\le5),\quad
S_j\cdot S_{j+1}=1\ (0\le j\le4),
\]
a pozostałe przecięcia różnych sfer są zerowe.
\end{definicja}

\begin{center}
\begin{tikzpicture}[x=1.23cm,y=1cm,line cap=round]
 \draw[tgTeal,line width=2.1pt] (0,0)--(5,0);
 \foreach \i in {0,...,5}{
   \ifnum\i=0
     \node[circle,draw=tgCoral,line width=1.7pt,fill=tgCoral!15,
       minimum size=10mm] (v\i) at (\i,0) {$-9$};
   \else
     \node[circle,draw=tgBlue,line width=1.4pt,fill=tgBlue!12,
       minimum size=9mm] (v\i) at (\i,0) {$-2$};
   \fi
   \node[font=\scriptsize,text=tgNavy] at (\i,-.74) {$S_{\i}$};
 }
 \node[font=\footnotesize,text=tgNavy] at (2.5,1.00)
   {graf ważony konfiguracji Parka $C_7$};
\end{tikzpicture}
\end{center}

\begin{twierdzenie}[Istnienie konfiguracji Parka; wynik zewnętrzny]
\label{thm:tg-park-embedding}
Istnieje wybór włókien osobliwych eliptycznego rozwarstwienia
$E(1)\to\mathbb{CP}^{1}$ oraz symplektyczne zanurzenie
$\iota\colon C_7\hookrightarrow R$ takie, że pięć sfer
$(-2)$ pochodzi z liniowego podgrafu osobliwego włókna
typu $\widetilde E_6$, a sfera $(-9)$ reprezentuje klasę
$s$ z przykładu \ref{ex:tg-park-S}. Ponadto dla
$V:=R\setminus\operatorname{int}\iota(C_7)$ zachodzi
$\pi_1(V)=1$.
\end{twierdzenie}

\begin{proof}[Konstrukcja i wskazanie nietrywialnego kroku]
W pracy J.~Parka, \emph{Simply connected symplectic
$4$-manifolds with $b_2^+=1$ and $c_1^2=2$},
Invent. Math. 159 (2005), 657--667, twierdzenie to
wynika z konstrukcji konfiguracji w Proposition~3.1 i
obliczenia grupy fundamentalnej w dowodzie Theorem~3.1.
Wybiera się cztery osobliwe włókna typu \emph{fishtail},
z których każde jest zanurzoną sferą z jednym dodatnim
punktem podwójnym, i rozdmuchuje ten punkt. Przekształcenie
właściwe takiego włókna ma klasę $f-2e_j$ i
samoprzecięcie $-4$; przecina sekcję $e_9$ dodatnio
w jednym punkcie. Po rozwiązaniu czterech przecięć z sekcją
powstaje jedna symplektyczna sfera w klasie
$s=4f+e_9-2\sum_{j=10}^{13}e_j$, której kwadrat
wynosi $-9$.

Dodatkowo włókno typu $\widetilde E_6$ zawiera łańcuch
pięciu sfer o samoprzecięciach $-2$. Ich przecięcia
z nową sferą dają dokładnie pokazany graf. Istnienie
odpowiednio dobranych włókien i ich symplektyczne
zanurzenie wymagają geometrii osobliwych rozwarstwień
eliptycznych, dlatego tego fragmentu nie zastępujemy
samym rachunkiem macierzy.

W dowodzie prostej spójności dopełnienia Park wskazuje
meridian konfiguracji reprezentujący generator
$\pi_1(Y_7)\cong\mathbb Z/49\mathbb Z$, który ogranicza
dysk w pozostałej części osobliwego włókna.
Wyjaśnijmy logiczny krok od tego faktu do tezy.
Ponieważ $R=V\cup_{Y_7} C_7$, obie części są spójne,
$\pi_1(R)=1$ i $\pi_1(C_7)=1$, twierdzenie van Kampena
mówi, że $\pi_1(V)$ jest normalnie generowana przez
obraz $\pi_1(Y_7)$. Skoro obraz generatora znika w
$\pi_1(V)$ dzięki wskazanemu dyskowi, to cały ten obraz
jest trywialny. Wynika stąd $\pi_1(V)=1$.
Istnienie wspomnianego dysku wymaga konkretnego układu
włókien z pracy Parka; ten geometryczny krok przyjmujemy
tutaj jako wynik zewnętrzny. Źródło: \url{https://doi.org/10.1007/s00222-004-0404-1}.
\end{proof}


\begin{lemat}[Jawne klasy pięciu sfer w $C_7$]
\label{lem:tg-park-S-classes}
W bazie $h,e_1,\ldots,e_{13}$ grupy $H_2(R;\mathbb Z)$
pięć sfer o kwadracie $-2$ można reprezentować klasami
\[
\begin{aligned}
 v_1&=e_4-e_7,& v_2&=e_1-e_4,&
 v_3&=h-e_1-e_2-e_3,\\
 v_4&=e_2-e_5,& v_5&=e_5-e_9.
\end{aligned}
\]
Ich kolejne przecięcia są równe $1$, a pozostałe przecięcia
różnych klas są zerowe. Klasa $s$ z przykładu
\ref{ex:tg-park-S} spełnia $s\cdot v_5=1$ i
$s\cdot v_j=0$ dla $1\le j\le4$.
\end{lemat}

\begin{proof}
Przez bezpośrednie podstawienie do diagonalnej formy przecięcia
obliczamy $v_j^2=-2$ oraz
$v_i\cdot v_{i+1}=1$ dla $1\le i\le4$.
Jeżeli $|i-j|>1$, odpowiednie klasy są ortogonalne.
Ponieważ $f\cdot v_j=0$, nowe klasy wyjątkowe
$e_{10},\ldots,e_{13}$ są ortogonalne do $v_j$, a
$e_9\cdot v_j=\delta_{5j}$, mamy
$s\cdot v_j=\delta_{5j}$.
Wynik nie zastępuje geometrycznego twierdzenia o realizacji
klas przez symplektyczne sfery; to pozostaje zewnętrznym
rezultatem Parka.
\end{proof}

\begin{uwaga}[Co rzeczywiście przedstawia diagram?]
Graf plumbingowy koduje samoprzecięcia i wzajemne
przecięcia sześciu sfer. \emph{Nie jest} pełnym diagramem
Kirby'ego rozmaitości $R$. Aby uzyskać pełny diagram,
trzeba również przedstawić uchwyty opisujące dopełnienie
$V$ oraz ich mapy przyczepienia. W tej iteracji rysujemy
poprawny lokalny diagram uchwytowy $C_7$, nie twierdząc,
że sam stanowi diagram całej konstrukcji Parka.
\end{uwaga}

\section{Macierz przecięcia i brzeg \texorpdfstring{$C_7$}{C7}}
\label{sec:tg-park-boundary}

\begin{stwierdzenie}[Macierz, określoność i wyznacznik]
\label{prop:tg-park-Q}
Macierz przecięcia $C_7$ w bazie $[S_0],\ldots,[S_5]$
ma postać
\[
Q_7=\begin{pmatrix}
-9&1&0&0&0&0\\
1&-2&1&0&0&0\\
0&1&-2&1&0&0\\
0&0&1&-2&1&0\\
0&0&0&1&-2&1\\
0&0&0&0&1&-2
\end{pmatrix}.
\]
Jest ujemnie określona, a jej wyznacznik wynosi $49$.
\end{stwierdzenie}

\begin{proof}
Wyraz diagonalny macierzy jest samoprzecięciem
odpowiedniej sfery; wyraz pozadiagonalny jest liczbą
przecięć dwóch sfer. Stąd od razu postać $Q_7$.

Niech $d_k$ będzie wyznacznikiem górnego lewego bloku
$k\times k$, z konwencją $d_0=1$. Wówczas
$d_1=-9$ i, przez rozwinięcie względem ostatniego wiersza,
$d_k=-2d_{k-1}-d_{k-2}$ dla $2\le k\le6$.
Indukcja daje $d_k=(-1)^k(8k+1)$: dla $k=0,1$
jest to prawda, a krok wynika z równości
$2(8k-7)-(8k-15)=8k+1$.
Dla $k=6$ otrzymujemy $\det Q_7=49$.

Ponadto $(-1)^k d_k=8k+1>0$ dla wszystkich $k$.
Z kryterium Sylvestera zastosowanego do $-Q_7$
wynika dodatnia określoność $-Q_7$, czyli ujemna
określoność $Q_7$.
\end{proof}

\begin{lemat}[Wszystkie współczynniki odwrotności macierzy Parka]
\label{lem:tg-park-Q-inverse-full}
Niech $Q_7$ będzie macierzą konfiguracji $C_7$
z Twierdzenia~\ref{prop:tg-park-Q}. Dla każdego
$0\le i,j\le5$ zachodzi
\[
(Q_7^{-1})_{ij}=
-\frac{(8\min\{i,j\}+1)(6-\max\{i,j\})}{49}.
\]
W szczególności diagonalne współczynniki formy
splecenia na meridianach sześciu składowych
plumbingu są w konwencji $\lambda=Q_7^{-1}$
równe modulo $\mathbb Z$
\[
-\frac1{49}(6,45,68,75,66,41).
\]
Każdy z sześciu meridianów reprezentuje generator
cyklicznej grupy $H_1(\partial C_7;\mathbb Z)$
rzędu $49$.
\end{lemat}

\begin{proof}
Niech $G=(g_{ij})$ będzie macierzą symetryczną
zadaną wzorem z tezy. Wykażemy bezpośrednio
$Q_7G=I_6$. Dla ustalonego $j$ i każdego
$1\le i\le4$, $i\ne j$, liczby $g_{ij}$
są liniowe względem $i$ po obu stronach
punktu $j$, zatem
$g_{i-1,j}-2g_{ij}+g_{i+1,j}=0$.
Dla $i=j$ ten sam wyraz wynosi
\[
\frac{-(8j-7)(6-j)+2(8j+1)(6-j)
-(8j+1)(5-j)}{49}=1.
\]
Na pierwszym wierszu sprawdzamy
$-9g_{0j}+g_{1j}=\delta_{0j}$,
a na ostatnim
$g_{4j}-2g_{5j}=\delta_{5j}$.
Wszystkie sześć wierszy iloczynu
$Q_7G$ daje więc macierz jednostkową.
Ponieważ $Q_7$ jest odwracalna,
$G=Q_7^{-1}$.

Podstawiając $i=j$ dla $i=0,\ldots,5$,
otrzymujemy kolejno wskazane liczniki
$6,45,68,75,66,41$. Każdy jest
względnie pierwszy z $49$.
Jeśli meridian $\mu_i$ miałby rząd
$n<49$ w $H_1(\partial C_7)$,
wówczas z biliniowości formy splecenia
$n\lambda(\mu_i,\mu_i)=0$ w
$\mathbb Q/\mathbb Z$. Ułamek o
mianowniku $49$ i liczniku
względnie pierwszym z $49$
wymusza $49\mid n$, sprzeczność.
Zatem wszystkie meridiany są generatorami.
\end{proof}

\begin{stwierdzenie}[Homologia brzegu]
\label{prop:tg-park-h1-boundary}
Dla $Y_7:=\partial C_7$ zachodzi
$H_1(Y_7;\mathbb Z)\cong\mathbb Z/49\mathbb Z$.
\end{stwierdzenie}

\begin{proof}
Ponieważ $C_7$ powstała z jednej $0$-komórki i sześciu
uchwytów indeksu dwa, jest po prostu spójna, a
$H_1(C_7)=0$, $H_2(C_7)\cong\mathbb Z^6$. Z ciągu
dokładnego pary $(C_7,Y_7)$ i dualności Poincarégo--Lefschetza
otrzymujemy
\[
H_2(C_7;\mathbb Z)\xrightarrow{Q_7}
H_2(C_7,Y_7;\mathbb Z)\longrightarrow
H_1(Y_7;\mathbb Z)\longrightarrow0.
\]
Zatem $H_1(Y_7)\cong\operatorname{coker}Q_7$.
Aby ustalić nie tylko rząd, lecz także strukturę grupy,
niech $x_0,\ldots,x_5$ odpowiadają generatorom prezentacji
przez kolumny $Q_7$. Ostatnie pięć relacji daje kolejno
$x_4=2x_5$, $x_3=3x_5$, $x_2=4x_5$,
$x_1=5x_5$ i $x_0=6x_5$.
Pierwsza relacja $-9x_0+x_1=0$ daje
$(-54+5)x_5=-49x_5=0$.
Wszystkie generatory są więc wielokrotnościami $x_5$,
a rząd grupy jest $|\det Q_7|=49$. Stąd grupa jest
cykliczna i ma podaną postać.
\end{proof}

\begin{stwierdzenie}[Wypełnienie soczewkowe]
\label{prop:tg-park-lens}
W konwencji, w której ujemny łańcuch $(-a_1,\ldots,-a_r)$
przedstawia chirurgię o współczynniku
$- [a_1,\ldots,a_r]^-$ na węźle trywialnym,
brzeg $C_7$ jest zorientowanie dyfeomorficzny z
$L(49,-6)$, gdzie
$[a_1,\ldots,a_r]^-=a_1-1/(a_2-\cdots-1/a_r)$.
\end{stwierdzenie}

\begin{proof}
Pięć końcowych dwójek redukuje się do ułamka $6/5$.
Stąd $[9,2,2,2,2,2]^-=9-5/6=49/6$.
Kolejne zastosowanie operacji \emph{slam-dunk}
(rachunku chirurgicznego torusów) zastępuje sześć
sprzężonych składowych pojedynczą wymierną chirurgią
o współczynniku $-49/6$. W konwencji chirurgii Dehna
przyjętej w skrypcie jest to przestrzeń soczewkowa
$L(49,-6)$. Samo obliczenie wyznacznika $49$ nie
wyznaczałoby parametru $-6$ ani orientowanego typu
rozmaitości, dlatego potrzebny jest rachunek ułamka.
\end{proof}

\begin{przyklad}[Forma splecenia: generator i orientacja]
Klasa meridianu przy składowej o ramowaniu $-9$
jest generatorem $H_1(Y_7)$. Rzeczywiście, lewy górny
minor macierzy $Q_7$ ma wyznacznik $-6$ (jest to macierz
łańcucha pięciu sfer o wagach $-2$). Wzór na odwrotność
macierzy daje więc $(Q_7^{-1})_{00}=-6/49$.
Ponieważ $6$ i $49$ są względnie pierwsze, element
takiej formy ma rząd $49$ w $\mathbb Q/\mathbb Z$;
stąd sama klasa meridianu ma rząd $49$.
Forma splecenia brzegu, z konwencją orientacji zgodną
z równaniem $\lambda_{Y_7}=Q_7^{-1}\pmod{\mathbb Z}$,
spełnia $\lambda_{Y_7}([\mu_0],[\mu_0])=-6/49$
w $\mathbb Q/\mathbb Z$. Przy odwróceniu orientacji
brzegu znak formy ulega zmianie. W literaturze spotyka
się też konwencję $-Q_7^{-1}$; nie należy mieszać
jej ze znakiem współczynnika chirurgii.
\end{przyklad}

\begin{center}
\begin{tikzpicture}[x=1.08cm,y=.95cm,>=Stealth,line cap=round]
  \foreach \i in {0,...,5}{
    \ifnum\i=0
      \draw[tgCoral,line width=1.9pt] (\i,0) circle (.32);
      \node[tgCoral,font=\scriptsize] at (\i,-.68) {$-9$};
    \else
      \draw[tgBlue,line width=1.9pt] (\i,0) circle (.32);
      \node[tgBlue,font=\scriptsize] at (\i,-.68) {$-2$};
    \fi
  }
  \foreach \i [evaluate=\i as \j using int(\i+1)] in {0,...,4}{
    \draw[tgTeal,line width=1.65pt] (\i+.32,.23)
          .. controls (\i+.52,.74) and (\j-.52,.74) .. (\j-.32,.23);
    \draw[tgTeal,line width=1.65pt] (\i+.32,-.23)
          .. controls (\i+.52,-.74) and (\j-.52,-.74) .. (\j-.32,-.23);
  }
  \node[font=\footnotesize,text=tgNavy] at (2.5,1.28)
       {Schemat przyległości składowych ramowania};
  \node[font=\footnotesize,text=tgNavy] at (2.5,-1.14)
       {Współczynnik ramowania jest podany pod składową};
\end{tikzpicture}

{\footnotesize Poglądowy schemat ramowań i sprzężeń: liczby pod okręgami
mają oznaczać ramowania, a połączenie sąsiadów ma
reprezentować liczbę splecenia $1$. Nie jest to regularny
diagram węzłowy z określonymi przejściami nad i pod;
ścisły diagram Kirby'ego wymaga narysowania tych skrzyżowań.}
\end{center}

\begin{lemat}[Metabolizer formy splecenia brzegu Parka]
\label{lem:tg-park-metabolizer}
Niech $Y_7=\partial C_7$ i niech $\mu_0$
oznacza meridian pierwszej składowej
ramowanego splotu o wadze $-9$.
W konwencji $\lambda_{Y_7}=Q_7^{-1}$
na grupie $H_1(Y_7;\mathbb Z)$ zachodzi
\[
H_1(Y_7;\mathbb Z)=\langle[\mu_0]\rangle
\cong\mathbb Z/49\mathbb Z,\qquad
\lambda_{Y_7}([\mu_0],[\mu_0])
=-\frac6{49}\pmod{\mathbb Z}.
\]
Podgrupa $M:=7\langle[\mu_0]\rangle$
rzędu $7$ jest równa swojemu dopełnieniu
ortogonalnemu $M^\perp$. W standardowym
wypełnieniu Cassona--Harera $B_7$
z Twierdzenia~\ref{thm:tg-park-B7}
jest ona jądrem surjekcji
$H_1(Y_7;\mathbb Z)\to H_1(B_7;\mathbb Z)$.
\end{lemat}

\begin{proof}
Z wcześniejszych obliczeń wiemy, że
$H_1(Y_7)\cong\mathbb Z/49\mathbb Z$.
Ponieważ $(Q_7^{-1})_{00}=-6/49$,
meridian $\mu_0$ ma rząd co najmniej $49$
(jeśli $n[\mu_0]=0$, to
$0=n\lambda([\mu_0],[\mu_0])=-6n/49$
w $\mathbb Q/\mathbb Z$ i $49\mid n$).
Zatem jest generatorem.

Element $7[\mu_0]$ ma rząd $7$ oraz
$\lambda(7[\mu_0],7[\mu_0])
=-6=0$ w $\mathbb Q/\mathbb Z$.
Dla $x[\mu_0]\in H_1(Y_7)$
warunek ortogonalności do $M$
przyjmuje postać
$-6x/7=0$ w $\mathbb Q/\mathbb Z$,
a więc $7\mid x$. Wynika stąd $M^\perp=M$.

Surjekcja grup fundamentalnych z Twierdzenia
\ref{thm:tg-park-B7} daje surjekcję
$H_1(Y_7)\to H_1(B_7)\cong\mathbb Z/7\mathbb Z$.
Jej jądro ma rząd $49/7=7$.
Cykliczna grupa rzędu $49$ ma dokładnie
jedną podgrupę rzędu $7$, mianowicie $M$.
\end{proof}

\section{Zastąpienie \texorpdfstring{$C_7$}{C7} wymierną kulą homologii}
\label{sec:tg-park-blowdown}

\begin{twierdzenie}[Casson--Harer; wynik zewnętrzny]
\label{thm:tg-park-B7}
Istnieje gładka, zwarta, orientowana czterorozmaitość
$B_7$ o brzegu $\partial B_7\cong L(49,-6)$, dla której
$H_*(B_7;\mathbb Q)\cong H_*(D^4;\mathbb Q)$ oraz
$\pi_1(B_7)\cong\mathbb Z/7\mathbb Z$.
Naturalne odwzorowanie
$\pi_1(\partial B_7)\to\pi_1(B_7)$ jest surjektywne.
\end{twierdzenie}

\begin{definicja}[Rozmaitość $X_7$]
Niech $V:=R\setminus\operatorname{int}\iota(C_7)$,
gdzie $\iota$ pochodzi z twierdzenia
\ref{thm:tg-park-embedding}. Wybierzmy
dyfeomorfizm $\partial B_7\to\partial V$ odwracający
orientacje indukowane na obu brzegach.
Definiujemy
\[
X_7:=V\cup_{\partial}B_7.
\]
Jest to \emph{wymierna kontrakcja konfiguracji $C_7$}
w powierzchni $R$.
\end{definicja}

\begin{center}
\begin{tikzpicture}[>=Stealth,line join=round,line cap=round]
  \node[draw=tgBlue,rounded corners=5pt,fill=tgBlue!9,
    minimum width=3.15cm,minimum height=1cm,align=center] (r)
     at (-3.3,1.08) {$R=V\cup C_7$\\$e=16,\ \sigma=-12$};
  \node[draw=tgCoral,rounded corners=5pt,fill=tgCoral!10,
    minimum width=2.8cm,minimum height=1cm,align=center] (c)
     at (1.85,1.08) {$C_7$\\$e=7,\ \sigma=-6$};
  \node[draw=tgTeal,rounded corners=5pt,fill=tgTeal!10,
    minimum width=2.8cm,minimum height=1cm,align=center] (b)
     at (1.85,-1.18) {$B_7$\\$e=1,\ \sigma=0$};
  \node[draw=tgPurple,rounded corners=5pt,fill=tgPurple!10,
    minimum width=3.0cm,minimum height=1cm,align=center] (x)
     at (-3.3,-1.18) {$X_7=V\cup B_7$\\$e=10,\ \sigma=-6$};
  \draw[-{Stealth[length=3mm]},line width=1.6pt,tgCoral]
        (r)-- node[above,font=\scriptsize] {usuń $C_7$} (c);
  \draw[-{Stealth[length=3mm]},line width=1.6pt,tgTeal]
        (c)-- node[right,font=\scriptsize,align=left]
        {zastąp przez\\$B_7$} (b);
  \draw[-{Stealth[length=3mm]},line width=1.6pt,tgPurple]
        (b)-- node[below,font=\scriptsize] {sklej z $V$} (x);
\end{tikzpicture}

{\footnotesize Schemat wymiernej kontrakcji. Równość brzegów umożliwia
operację wycięcia i wklejenia; sam schemat nie dowodzi
zachowania prostej spójności.}
\end{center}

\begin{twierdzenie}[Grupa fundamentalna wyniku]
\label{thm:tg-park-pi1}
Rozmaitość $X_7$ jest po prostu spójna.
\end{twierdzenie}

\begin{proof}
Rozważmy rozkład $X_7=V\cup B_7$,
którego część wspólna jest kołnierzem spójnej
przestrzeni soczewkowej $Y_7$.
Twierdzenie van Kampena daje
\[
\pi_1(X_7)\cong
\pi_1(V)*_{\pi_1(Y_7)}\pi_1(B_7).
\]
Z twierdzenia \ref{thm:tg-park-embedding} mamy
$\pi_1(V)=1$. Zatem prawa strona jest ilorazem
$\pi_1(B_7)$ przez domknięcie normalne obrazu
$\pi_1(Y_7)$. Z twierdzenia \ref{thm:tg-park-B7}
odwzorowanie $\pi_1(Y_7)\to\pi_1(B_7)$ jest
surjektywne, a więc iloraz jest trywialny.
\end{proof}

\begin{twierdzenie}[Charakterystyki liczbowe]
\label{thm:tg-park-numbers}
Zachodzą
\[
e(X_7)=10,\quad \sigma(X_7)=-6,\quad
b_1(X_7)=b_3(X_7)=0,\quad
b_2^+(X_7)=1,\quad b_2^-(X_7)=7.
\]
Ponadto $H_2(X_7;\mathbb Z)\cong\mathbb Z^8$.
\end{twierdzenie}

\begin{proof}
$C_7$ powstaje z jednego uchwytu indeksu zero
i sześciu uchwytów indeksu dwa. Wobec tego
$e(C_7)=1+6=7$; nie jest to $8$, ponieważ
liczba wierzchołków grafu liczy uchwyty indeksu
dwa, a nie wszystkie komórki. Jej forma jest
ujemnie określona rzędu sześć, więc
$\sigma(C_7)=-6$.
Wymierna kula homologii ma $e(B_7)=1$ i
$\sigma(B_7)=0$.
Ponieważ brzeg $Y_7$ ma zerową charakterystykę
Eulera, z addytywności otrzymujemy
\[
e(X_7)=e(R)-e(C_7)+e(B_7)=16-7+1=10.
\]
Z addytywności sygnatury (twierdzenie Novikova)
przy rozkładzie wzdłuż wymiernej sfery homologicznej
wynika $\sigma(X_7)=\sigma(R)-\sigma(C_7)=-12+6=-6$.
Prosta spójność daje $H_1(X_7;\mathbb Z)=0$, więc
$b_1=0$. Dualność Poincarégo daje $b_3=0$.
Zatem $e=2+b_2$, stąd $b_2(X_7)=8$.
Rozwiązanie równań $b_2^++b_2^-=8$ i
$b_2^+-b_2^-=-6$ prowadzi do $(b_2^+,b_2^-)=(1,7)$.
Wreszcie uniwersalne twierdzenie o współczynnikach
i dualność Poincarégo dają brak torsji w
$H_2(X_7;\mathbb Z)$ dla po prostu spójnej zamkniętej
orientowanej czterorozmaitości. Skoro rząd jest osiem,
$H_2(X_7;\mathbb Z)\cong\mathbb Z^8$.
\end{proof}

\begin{przyklad}[Kontrola wzoru $c_1^2$]
Twierdzenie Symington o symplektycznej wymiernej kontrakcji
zapewnia, że $X_7$ ma strukturę symplektyczną, ponieważ
konfiguracja Parka składa się z symplektycznych sfer
przecinających się dodatnio. Dla zgodnej struktury prawie
zespolonej na takiej zamkniętej czterorozmaitości zachodzi
$c_1^2=2e+3\sigma$. Dla $X_7$ otrzymujemy
$c_1^2(X_7)=20-18=2$.
Zgadza się to z twierdzeniem Parka. Sama równość
$c_1^2=2$ nie rozróżnia jednak struktur gładkich:
tę samą liczbę ma standardowa racjonalna powierzchnia
$\mathbb{CP}^{2}\#7\overline{\mathbb{CP}}^{\,2}$.
\end{przyklad}

\section{Forma przecięcia i typ homeomorfizmu}
\label{sec:tg-park-homeo}

\begin{lemat}[Nieparzystość formy przecięcia]
\label{lem:tg-park-odd}
Forma przecięcia $X_7$ jest nieparzysta.
\end{lemat}

\begin{proof}
Gdyby forma przecięcia $Q_{X_7}$ była parzysta,
rozmaitość $X_7$ byłaby spin: dla po prostu spójnej
orientowanej czterorozmaitości warunek
$w_2(X_7)=0$ jest równoważny parzystości formy.
Twierdzenie Rochlina mówi, że sygnatura zamkniętej
gładkiej czterorozmaitości spin jest podzielna przez $16$.
Tymczasem $\sigma(X_7)=-6$, sprzeczność.
Wynika stąd nieparzystość formy.
\end{proof}

\begin{twierdzenie}[Topologiczny typ $X_7$]
\label{thm:tg-park-homeo}
Istnieje zachowujący orientację homeomorfizm
\[
X_7\cong_{\mathrm{TOP}}
\mathbb{CP}^{2}\#7\overline{\mathbb{CP}}^{\,2}.
\]
\end{twierdzenie}

\begin{proof}
Dla każdej zamkniętej, spójnej, orientowanej, po prostu
spójnej czterorozmaitości forma przecięcia na
$H_2/\operatorname{Tor}$ jest unimodularna.
Twierdzenie \ref{thm:tg-park-numbers} i lemat
\ref{lem:tg-park-odd} dają więc nieparzystą,
nieokreśloną formę unimodularną o sygnaturze $-6$
i rzędzie osiem. Klasyfikacja nieparzystych
nieokreślonych form unimodularnych nad $\mathbb Z$
implikuje izometrię z formą
$\langle1\rangle\oplus7\langle-1\rangle$.

Aby przejść od formy do rozmaitości, stosujemy twierdzenie
klasyfikacyjne Freedmana dla zamkniętych, po prostu
spójnych topologicznych czterorozmaitości.
Obie rozmaitości są wygładzalne, zatem ich niezmiennik
Kirby'ego--Siebenmanna jest zerowy. Dla ustalonej
nieparzystej formy i tego niezmiennika dostajemy
homeomorfizm zachowujący orientację.
Twierdzenie Freedmana przyjmujemy tu jako wynik
zewnętrzny, natomiast identyfikację samej formy
wyprowadziliśmy z poprzednich rachunków.
\end{proof}

\begin{uwaga}[Czego jeszcze nie dowiedliśmy]
W dotychczasowych rachunkach uzyskaliśmy homeomorfizm,
nie zaś dyfeomorfizm. Istnienie formy przecięcia
$\langle1\rangle\oplus7\langle-1\rangle$ nie pozwala
rozpoznać gładkiej struktury. Należy porównać
niezmienniki Seiberga--Wittena w odpowiedniej komorze
przy $b_2^+=1$, wykorzystując symplektyczność $X_7$,
odpowiednie twierdzenie Taubesa--Li--Liu i nierówność
dodatniego parowania klasy kanonicznej z klasą formy
symplektycznej. W następnym rozdziale przeprowadzimy ten argument.
\end{uwaga}

\section{Kontrolny rachunek macierzy i granice metody}
\label{sec:tg-park-summary}

\begin{przyklad}[Rekurencja ułamka łańcuchowego]
Oznaczmy przez $r_j$ ujemny ułamek łańcuchowy
złożony z $j$ dwójek: $r_1=2$ oraz
$r_{j+1}=2-1/r_j$. Indukcja daje
$r_j=(j+1)/j$. Stąd dla pięciu dwójek
otrzymujemy $r_5=6/5$, a więc
$[9,2,2,2,2,2]^-=9-1/r_5=49/6$.
To drugi, niezależny od rachunku wyznacznika sposób
ustalenia liczby $49$ i parametru chirurgii $-6$.
\end{przyklad}

\begin{przyklad}[Bilans liczb Betti'ego]
Rozmaitość $R$ ma $b_2^-=13$. Wycięta konfiguracja
$C_7$ wnosi sześć ujemnych kierunków, a $B_7$ ma
zerową drugą homologię wymierną. Zatem $b_2^-$
spada do $7$, podczas gdy $b_2^+=1$ pozostaje
bez zmian. Zgadza się to z obliczeniem przez
$(e,\sigma)$: w obu metodach otrzymujemy $b_2=8$.
\end{przyklad}

\begin{uwaga}[Zakres dowodu a konstrukcja globalna]
Z rozdziału można samodzielnie odtworzyć wszystkie
rachunki na grafie $C_7$, obliczenie brzegu, zastosowanie
van Kampena \emph{przy znanej prostej spójności dopełnienia},
liczby Betti'ego oraz formę przecięcia i typ TOP.
Zaawansowanymi danymi wejściowymi pozostają
(1) istnienie konkretnego symplektycznego zanurzenia
$C_7\hookrightarrow R$ i dowód $\pi_1(V)=1$ z pracy Parka,
(2) istnienie wymiernej kuli homologii $B_7$ i jej
własności z konstrukcji Cassona--Harera,
(3) twierdzenie Freedmana.

W kolejnym rozdziale do tej listy dojdzie obliczenie
niezmienników Seiberga--Wittena i porównanie komór.
Żaden z tych wyników nie wynika z samej macierzy $Q_7$.
\end{uwaga}


% ========================================================================
% Iteracja 14. Umiescic bezposrednio po rozdziale o konstrukcji Parka.
% Konwencja: u_0=S_0 (sfera -9), u_1=S_5 z pracy Parka,
% ..., u_5=S_1 z pracy Parka. Tylko ponizszy dokument zawiera pelne
% doprowadzenie do nierownosci K_{X_7}.[omega_7]>0.
% ========================================================================
\begin{uwaga}[Geometria konfiguracji Parka]
Rachunek $Q_{C_7}$ oraz równości
$e(X_7)=10$, $\sigma(X_7)=-6$ nie konstruują same
sześciu symplektycznych sfer w rozdmuchanej powierzchni
eliptycznej. Ich konkretne klasy i istnienie zgodnej
formy symplektycznej są danymi geometrycznymi konstrukcji
Parka. Podobnie prosta spójność dopełnienia musi zostać
wykazana przez odpowiedni argument o meridianie,
a nie wywnioskowana wyłącznie z wyznacznika formy.
\end{uwaga}

